\documentclass[12pt, letterpaper]{article}

\usepackage[utf8]{inputenc}
\usepackage[T1]{fontenc}
\usepackage[spanish]{babel}
\usepackage{geometry}
\usepackage{amsmath, amssymb, amsthm}
\usepackage{booktabs}
\usepackage{array}
\usepackage{microtype}
\usepackage{setspace}
\usepackage{parskip}
\usepackage{titlesec}
\usepackage{fancyhdr}
\usepackage{enumitem}
\usepackage{bm}
\usepackage{mathtools}
\usepackage{tcolorbox}
\usepackage{hyperref}

\tcbuselibrary{skins, breakable}

\geometry{top=2.8cm, bottom=3.0cm, left=3.2cm, right=3.2cm, headheight=14pt}
\setstretch{1.20}

\pagestyle{fancy}
\fancyhf{}
\fancyhead[C]{\small\textsc{Econometría --- Gujarati --- Ejercicio 13.20}}
\fancyfoot[C]{\small\thepage}
\renewcommand{\headrulewidth}{0.4pt}
\renewcommand{\footrulewidth}{0pt}

\titleformat{\section}
  {\large\bfseries\scshape}{\thesection.}{0.6em}{}
  [\vspace{-4pt}\rule{\linewidth}{0.4pt}]
\titleformat{\subsection}{\normalsize\bfseries}{\thesubsection.}{0.5em}{}
\titleformat{\subsubsection}{\normalsize\itshape}{}{0em}{}

\newtcolorbox{clave}{
  colback=white, colframe=black,
  boxrule=0.5pt, arc=3pt,
  left=8pt, right=8pt, top=6pt, bottom=6pt,
  breakable
}

\newcommand{\E}{\operatorname{E}}
\newcommand{\Cov}{\operatorname{Cov}}
\newcommand{\Var}{\operatorname{Var}}

\begin{document}

\begin{center}
  {\LARGE\bfseries Observaciones Influyentes, Valores Atípicos,}\\[6pt]
  {\LARGE\bfseries Jerarquía Polinomial y Variables Centradas}\\[12pt]
  {\large\itshape Verdadero o Falso: cinco enunciados sobre diagnóstico}\\
  {\large\itshape y especificación en regresión lineal}\\[14pt]
  \rule{0.55\linewidth}{0.6pt}\\[8pt]
  {\normalsize Ejercicio~13.20 --- \textit{Econometría}, Gujarati \& Porter, 5.a ed.}\\[4pt]
  {\small\textsc{Análisis conceptual, teórico y algebraico}}
\end{center}

\vspace{10pt}

%------------------------------------------------------------------------------
\section{Enunciado a) Una observación puede ser influyente pero no atípica}
%------------------------------------------------------------------------------

\subsection*{\textit{Veredicto: \textbf{Verdadero}}}

\subsection{Distinción conceptual}

En el análisis de regresión se distinguen dos conceptos que habitualmente
se confunden:

Un \textbf{valor atípico} (\textit{outlier}) es una observación cuyo
\emph{residual} $\hat{u}_i$ es grande en valor absoluto, es decir, un
punto que se aleja verticalmente de la línea de regresión estimada.

Una \textbf{observación influyente} es aquella que, al ser eliminada,
produce un cambio sustancial en los coeficientes estimados
$\hat{\boldsymbol{\beta}}$. La influencia depende no solo del residual
sino también de la posición del punto en el espacio de los regresores.

\subsection{Argumento teórico: el papel del apalancamiento}

La medida de \textbf{apalancamiento} (\textit{leverage}) de la
observación $i$ es el $i$-ésimo elemento diagonal de la matriz sombrero
(\textit{hat matrix}) $\mathbf{H} = \mathbf{X}(\mathbf{X}'\mathbf{X})^{-1}\mathbf{X}'$:
\begin{equation}
  h_{ii} = \mathbf{x}_i'(\mathbf{X}'\mathbf{X})^{-1}\mathbf{x}_i
  \in \left[\frac{1}{n},\, 1\right],
  \label{eq:leverage}
\end{equation}
donde $\mathbf{x}_i$ es el vector de regresores de la observación $i$.
Un valor $h_{ii}$ elevado indica que la observación está muy alejada de
la media muestral de los regresores $\bar{\mathbf{X}}$, lo que le otorga
un gran poder para ``atraer'' la recta de regresión hacia sí misma.

\subsection{Mecanismo}

Cuando $h_{ii}$ es muy grande, la línea de regresión MCO se ajusta
casi perfectamente a esa observación para minimizar la suma de cuadrados,
produciendo un residual $\hat{u}_i \approx 0$. El punto no se detecta
como atípico al examinar los residuales, pero al eliminarlo la pendiente
estimada cambia drásticamente. Se dice entonces que es \emph{influyente
pero no atípico}.

\subsection{Argumento algebraico}

El cambio en los estimadores MCO al omitir la observación $t$ es:
\begin{equation}
  \hat{\boldsymbol{\beta}}_{(t)} - \hat{\boldsymbol{\beta}}
  = -\frac{1}{1 - h_{tt}}
    (\mathbf{X}'\mathbf{X})^{-1}\mathbf{x}_t\,\hat{u}_t,
  \label{eq:cambio_coef}
\end{equation}
donde $h_{tt}$ es el apalancamiento de la observación $t$. El factor
$1/(1 - h_{tt})$ actúa como un amplificador: a medida que $h_{tt} \to 1$,
cualquier residual, por pequeño que sea, genera un cambio enorme en
los coeficientes. Así, una observación con $h_{tt}$ cercano a 1 y
$\hat{u}_t \approx 0$ no se detecta como atípica, pero sí resulta
influyente.

\begin{clave}
  \textbf{Conclusión a):} Una observación puede tener alto apalancamiento
  ($h_{ii}$ cercano a 1), residual pequeño (no atípica) y, sin embargo,
  modificar sustancialmente los coeficientes al ser eliminada (influyente).
  La influyencia y la condición de atípica son propiedades independientes.
\end{clave}

%------------------------------------------------------------------------------
\section{Enunciado b) Una observación puede ser atípica pero no influyente}
%------------------------------------------------------------------------------

\subsection*{\textit{Veredicto: \textbf{Verdadero}}}

\subsection{Argumento conceptual}

Una observación con residual grande $|\hat{u}_t|$ (atípica) puede tener
bajo apalancamiento $h_{tt}$ si se ubica cerca de la media de los
regresores $\bar{\mathbf{X}}$. En ese caso, la observación se aleja
verticalmente de la línea de regresión, pero su posición ``centrada''
en el espacio de las $X$ le impide modificar significativamente la
dirección de esa línea.

\subsection{Argumento algebraico}

Retomando~\eqref{eq:cambio_coef}, si $h_{tt}$ es pequeño (próximo a
$1/n$), el factor amplificador $1/(1 - h_{tt}) \approx n/(n-1) \approx 1$
resulta modesto. Incluso con un residual $|\hat{u}_t|$ grande, el cambio
en los coeficientes $\hat{\boldsymbol{\beta}}_{(t)} - \hat{\boldsymbol{\beta}}$
será pequeño. La observación es un \emph{outlier de bajo apalancamiento}:
distorsiona el ajuste local pero no la dirección global de la recta.

\begin{clave}
  \textbf{Conclusión b):} Una observación puede tener residual grande
  (atípica) y, a la vez, apalancamiento bajo (no influyente). La
  condición de valor atípico se define por el residual; la influencia,
  por la combinación de residual y apalancamiento.
\end{clave}

%------------------------------------------------------------------------------
\section{Enunciado c) Una observación puede ser atípica e influyente}
%------------------------------------------------------------------------------

\subsection*{\textit{Veredicto: \textbf{Verdadero}}}

\subsection{Argumento teórico}

Este caso ocurre cuando una observación combina \emph{simultáneamente}
alto apalancamiento ($h_{tt}$ grande) y un comportamiento inconsistente
con el patrón del resto de los datos. El punto está lejos de la media
de los regresores (alto leverage) y, además, no sigue la relación lineal
prevaleciente. La línea MCO intenta reducir el residual desplazándose
hacia el punto, pero no puede eliminarlo por completo, de modo que el
residual sigue siendo grande.

\subsection{Diagnóstico operativo: la distancia de Cook}

Para detectar este caso se utiliza la \textbf{distancia de Cook}:
\begin{equation}
  D_i = \frac{(\hat{\boldsymbol{\beta}}_{(i)} - \hat{\boldsymbol{\beta}})'
              \mathbf{X}'\mathbf{X}
              (\hat{\boldsymbol{\beta}}_{(i)} - \hat{\boldsymbol{\beta}})}
             {k\,\hat{\sigma}^2},
  \label{eq:cook}
\end{equation}
donde $k$ es el número de parámetros. Una $D_i$ grande señala que la
observación $i$ es \emph{a la vez} influyente y atípica, pues requiere
tanto un residual sustancial como un apalancamiento elevado para producir
el desplazamiento $\hat{\boldsymbol{\beta}}_{(i)} - \hat{\boldsymbol{\beta}}$
notable.

\begin{clave}
  \textbf{Conclusión c):} Los enunciados a), b) y c) son todos verdaderos
  y, tomados en conjunto, muestran que atipicidad e influencia son
  propiedades que pueden coexistir o no. El diagnóstico completo requiere
  examinar tanto los residuales como el apalancamiento y la distancia
  de Cook, no solo uno de estos indicadores.
\end{clave}

%------------------------------------------------------------------------------
\section{Enunciado d) Si $\hat{\beta}_3$ es significativo en el modelo cuadrático,
         se conserva $X_i$ aunque $\hat{\beta}_2$ sea insignificante}
%------------------------------------------------------------------------------

\subsection*{\textit{Veredicto: \textbf{Verdadero}}}

\subsection{El principio de jerarquía en modelos polinomiales}

Considérese el modelo cuadrático:
\begin{equation}
  Y_i = \beta_1 + \beta_2 X_i + \beta_3 X_i^2 + u_i.
  \label{eq:cuadratico}
\end{equation}
Si $\hat{\beta}_3$ resulta estadísticamente significativo, el modelo
contiene un término de \emph{orden superior} validado por los datos.
El \textbf{principio de jerarquía} establece que, cuando un término
polinomial de orden $p$ se incluye en el modelo, todos los términos
de orden inferior ($1, X, X^2, \ldots, X^{p-1}$) deben conservarse,
independientemente de su significancia individual.

\subsection{Argumento algebraico: el efecto marginal}

El efecto parcial de $X_i$ sobre $Y_i$ en el modelo~\eqref{eq:cuadratico}
es:
\begin{equation}
  \frac{\partial \hat{Y}_i}{\partial X_i}
  = \hat{\beta}_2 + 2\hat{\beta}_3 X_i.
  \label{eq:efecto_marginal}
\end{equation}
Suprimir $\beta_2$ del modelo equivale a imponer la restricción
$\beta_2 = 0$, lo que fuerza a que la pendiente sea exactamente cero
en el punto $X_i = 0$. Esta restricción no tiene justificación
económica general y distorsiona el efecto marginal para todos los
valores de $X_i$, incluyendo los del interior de la muestra.

\subsection{Argumento sobre la reparametrización}

El término $\hat{\beta}_2$ siendo individualmente insignificante no
implica que $X_i$ sea irrelevante; puede reflejar simplemente
multicolinealidad entre $X_i$ y $X_i^2$. Si se elimina $X_i$ y se
estima el modelo restringido $Y_i = \beta_1 + \beta_3 X_i^2 + u_i$,
los estimadores de $\beta_1$ y $\beta_3$ cambian de interpretación y
absorben parte del efecto lineal de $X_i$, produciendo sesgos en toda
la especificación.

\begin{clave}
  \textbf{Conclusión d):} En un modelo polinomial, la significancia de
  $\hat{\beta}_3$ basta para mantener $X_i$ en el modelo aunque
  $\hat{\beta}_2$ sea individualmente insignificante. Eliminar el
  término lineal impone una restricción funcional arbitraria sobre el
  efecto marginal y sesga los demás coeficientes.
\end{clave}

%------------------------------------------------------------------------------
\section{Enunciado e) La línea de regresión es la misma en niveles y en desviaciones}
%------------------------------------------------------------------------------

\subsection*{\textit{Veredicto: \textbf{Verdadero}}}

\subsection{Planteamiento}

Se comparan dos modelos. El \textbf{modelo en niveles}:
\begin{equation}
  Y_i = \beta_1 + \beta_2 X_{2i} + \beta_3 X_{3i} + u_i,
  \label{eq:niveles}
\end{equation}
y el \textbf{modelo en desviaciones respecto a la media}, con
$x_{2i} = X_{2i} - \bar{X}_2$ y $x_{3i} = X_{3i} - \bar{X}_3$:
\begin{equation}
  Y_i = \alpha_1 + \beta_2 x_{2i} + \beta_3 x_{3i} + u_i.
  \label{eq:desviaciones}
\end{equation}

\subsection{Demostración algebraica}

En el modelo en desviaciones~\eqref{eq:desviaciones}, los estimadores
MCO de las pendientes son idénticos a los de~\eqref{eq:niveles}, pues
centrar los regresores es una transformación lineal que no altera las
covarianzas ni las varianzas relativas entre variables. Formalmente:
\begin{equation}
  \hat{\beta}_j^{\text{desv}}
  = \frac{\sum_i x_{ji}\,(Y_i - \bar{Y})}{\sum_i x_{ji}^2}
  = \frac{\sum_i (X_{ji} - \bar{X}_j)(Y_i - \bar{Y})}
         {\sum_i (X_{ji} - \bar{X}_j)^2}
  = \hat{\beta}_j^{\text{niv}}, \quad j = 2, 3.
  \label{eq:igualdad_pendientes}
\end{equation}
Para el intercepto, la condición de primer orden del MCO en el modelo
en desviaciones implica que la recta pasa por $(\bar{x}_2, \bar{x}_3, \bar{Y})
= (0, 0, \bar{Y})$, de modo que:
\begin{equation}
  \hat{\alpha}_1 = \bar{Y}.
  \label{eq:intercepto_desv}
\end{equation}
El intercepto del modelo en niveles se obtiene de la misma condición de
primer orden:
\begin{equation}
  \hat{\beta}_1 = \bar{Y} - \hat{\beta}_2\bar{X}_2 - \hat{\beta}_3\bar{X}_3.
  \label{eq:intercepto_niv}
\end{equation}
Sustituyendo~\eqref{eq:intercepto_niv} y las pendientes iguales en los
valores ajustados del modelo en niveles:
\begin{align}
  \hat{Y}_i^{\text{niv}}
  &= \hat{\beta}_1 + \hat{\beta}_2 X_{2i} + \hat{\beta}_3 X_{3i} \notag\\
  &= (\bar{Y} - \hat{\beta}_2\bar{X}_2 - \hat{\beta}_3\bar{X}_3)
     + \hat{\beta}_2 X_{2i} + \hat{\beta}_3 X_{3i} \notag\\
  &= \bar{Y} + \hat{\beta}_2(X_{2i} - \bar{X}_2)
             + \hat{\beta}_3(X_{3i} - \bar{X}_3) \notag\\
  &= \hat{\alpha}_1 + \hat{\beta}_2 x_{2i} + \hat{\beta}_3 x_{3i}
   = \hat{Y}_i^{\text{desv}}.
  \label{eq:igualdad_yhat}
\end{align}
Los valores ajustados son numéricamente idénticos en ambas
especificaciones y, por tanto, los residuales también lo son:
\begin{equation}
  \hat{u}_i^{\text{niv}} = Y_i - \hat{Y}_i^{\text{niv}}
  = Y_i - \hat{Y}_i^{\text{desv}}
  = \hat{u}_i^{\text{desv}}.
  \label{eq:igualdad_residuales}
\end{equation}
Dado que los valores ajustados y los residuales son idénticos, el
hiperplano de regresión estimado es exactamente el mismo en ambas
especificaciones: solo el intercepto toma un valor numérico diferente
porque mide el nivel de $Y$ en un punto de referencia distinto
($\mathbf{X} = \mathbf{0}$ en niveles; $\mathbf{x} = \mathbf{0}$,
es decir, $\mathbf{X} = \bar{\mathbf{X}}$, en desviaciones).

\begin{clave}
  \textbf{Conclusión e):} Las dos especificaciones producen exactamente
  la misma línea (hiperplano) de regresión, los mismos valores ajustados
  y los mismos residuales. La diferencia es únicamente de reparametrización:
  el intercepto del modelo centrado ($\hat{\alpha}_1 = \bar{Y}$) es el
  valor ajustado en el punto de las medias, mientras que el intercepto
  del modelo en niveles es el valor ajustado en $\mathbf{X} = \mathbf{0}$.
  Las pendientes son numéricamente idénticas en ambos modelos.
\end{clave}

\end{document}
